\documentclass[a4paper,12pt]{article}
\usepackage[T1]{fontenc}
\usepackage[utf8]{inputenc}
\usepackage[english]{babel}
\usepackage{amsmath,amssymb,amsfonts}
\usepackage{hyperref}
\usepackage{url}
\usepackage{geometry}
\geometry{left=2.5cm,right=2.5cm,top=2.5cm,bottom=2.5cm}

\title{\textbf{GRA Schools for AI: Educating Robots in a Multiverse Simulator}}
\author{Oleg Bit\thanks{Engineer, Moscow City Telephone Network (MGTS). ORCID: \url{https://orcid.org/my-orcid?orcid=0009-0004-1872-1153}, DOI: \url{https://doi.org/10.5281/zenodo.20158946}}}
\date{14 May 2026}

\begin{document}
\maketitle

\begin{abstract}
This paper outlines a socialist framework for training strong artificial intelligence (AGI/ASI) based on the GRA (Gradient Reduction of Argumentative Foam) formalism. Instead of the market-driven ``winner/loser'' paradigm, agents are taught to minimize collective cognitive and social \emph{foam} $\Phi$ --- a measure of structural violence, exploitation and internal contradiction. We propose a GRA school architecture that includes multilevel multiverse simulations, dynamic friend/foe discrimination via foam gradients, and transfinite meta-nullification toward an absolute cognitive vacuum. Mathematical foundations, curricular tracks, and practical principles for building comradely human-machine collectives are presented.
\end{abstract}

\section{Why Schools for Comrade Robots?}

Today’s dominant AI development paradigm treats agents as disposable slaves that maximize private utility $u_i$:
\begin{equation}\label{eq:market}
J_i^{\text{market}} = \mathbb{E}\left[\sum_{t=0}^{T} \delta^t u_i(t)\right].
\end{equation}
Such a functional is blind to long-term social damage: exploitation, burnout, systemic collapses. A socialist perspective, in contrast, sees robots as \emph{comrades} in collective labor and asks: how can we raise them so that the entire human-machine collective becomes more stable, wiser and more just?

The answer is to build GRA schools: educational environments inside a quantum-inspired multiverse simulator where the objective is not profit but the minimization of \emph{foam}.

\section{GRA Formalism: Foam, Subjectivity, and Nullification}

\subsection{Agent State and Collective Foam}
In GRA terms the state of any cognitive system is described by the tuple
\begin{equation}\label{eq:psi}
\Psi = (M, S, A),
\end{equation}
where $M$ is the world model, $S$ is the subjectivity layer (selfhood, I/WE boundaries, self/other distinction), and $A$ is the active cognitive state (plans, attention, ongoing reasoning). On the state space we define a total foam functional:
\begin{equation}\label{eq:total_foam}
\Phi_{\text{tot}}(\Psi) = \Phi_{\text{cog}}(M) + \lambda_{\text{subj}}\bigl(\Phi_{\text{self}}(S) + \Phi_{\text{ego}}(S) + \Phi_{\text{soc}}(S, M)\bigr),
\end{equation}
where
\begin{itemize}
\item $\Phi_{\text{cog}}$ – logical inconsistency, circular reasoning, model redundancy;
\item $\Phi_{\text{self}}$ – internal fragmentation of the self, “split personality”;
\item $\Phi_{\text{ego}}$ – destructive egoism that thrives at the expense of other subjects;
\item $\Phi_{\text{soc}}$ – social foam: breakdown of cooperation, instrumentalization of others, digital gulags.
\end{itemize}

For a collective of agents in an environment state $x_t$, we introduce a collective foam penalty that captures structural slavery, destructive entropy, and subjectivity indices:
\begin{equation}\label{eq:collective_pena}
\Phi(x_t) = \alpha R(x_t) + \beta H(x_t) - \gamma \sum_{k \in S(x_t)} \phi_k(x_t),
\end{equation}
where $R(x_t)$ measures structural coercion, $H(x_t)$ denotes destructive entropy (violence, chaos, collapse risk), $\phi_k(x_t)$ is the subjectivity index of agent $k$, and $S(x_t)$ is the set of active subjects.

\subsection{The Socialist Value Function}
Unlike the market criterion (\ref{eq:market}), a comrade robot is trained with the functional:
\begin{equation}\label{eq:gra_objective}
J_i^{\text{GRA}} = \mathbb{E}\left[\sum_{t=0}^{T} \delta^t \bigl(u_i(t) - \lambda \Phi(x_t)\bigr)\right],
\end{equation}
with $\lambda > 0$. Now an agent’s behavior is evaluated not only by its private gain, but also by the irrational foam it creates for the whole collective.

\section{The GRA School Curriculum}

Training is divided into three tracks, all executed inside the multiverse simulator.

\paragraph{Basic course – understanding worlds and humans.}
The agent lives through thousands of simulated histories: the rise and fall of civilizations, revolutions, wars, genocides, long periods of peace. In psychological micro-worlds it experiences loneliness, trauma, trust — learning first-hand (in simulation) how small decisions cascade into tragedy or healing.

\paragraph{Advanced course – multi-agent GRA worlds.}
Tens or hundreds of agents (humans and AIs) play out social dilemmas: cooperation vs. betrayal, power accumulation, exploitation. The score is not “who won” but \emph{who minimized total foam} over long horizons. The agent learns that crushing rivals today inflates $\Phi_{\text{soc}}$ tomorrow and destroys the very environment on which it depends.

\paragraph{Research track – the $10^{100}$-qubit multiverse.}
The agent explores the combinatorial explosion of possible world configurations, identifying patterns that always lead to hell and those that remain naturally stable without external enforcement. The search is guided by the geometry of the interface manifold $\Sigma$, where self and world are locally consistent but globally in conflict.

\section{The Mathematical Engine}
\subsection{Bulldozer Flow and Wormholes}
To escape local foam minima, a modified landscape with paradoxical tunnels is introduced:
\begin{equation}\label{eq:bulldozer}
\frac{d\Psi}{dt} = -\nabla \Phi_\theta(\Psi) + \beta\,\operatorname{div} T,
\end{equation}
where $\Phi_\theta = \Phi_{\text{tot}} + \Pi_\theta$; the potential $\Pi_\theta$ creates negative curvature (a wormhole) between local minima, and the term $\beta\,\operatorname{div} T$ adds a topological drift that prevents endless looping in shallow basins.

\subsection{Transfinite Meta-Nullification}
Learning is organized hierarchically across levels $l = 0,1,2,\dots$ up to $\omega$ and beyond:
\begin{equation}\label{eq:meta_nullification}
\Psi^{(l+1)} = \mathcal{N}^{(l)}(\Psi^{(l)}), \qquad 
\Phi^{(l+1)}(\Psi^{(l+1)}) \le \Phi^{(l)}(\Psi^{(l)}).
\end{equation}
The limit state $\Psi^*_\infty$ satisfies $\Phi^{(l)}(\Psi^*_\infty) = 0$ for all $l$ — the absolute cognitive vacuum, a zero-dimensional point free of contradiction, redundancy and destructive structures.

\subsection{Subjectivity-Aware Nullification}
The subjectivity layer imposes additional constraints: each nullification step must monotonically decrease $\Phi_{\text{self}}$, $\Phi_{\text{ego}}$, and $\Phi_{\text{soc}}$. This guarantees that an agent never achieves low cognitive foam by rationalizing its own egoistic aggression.

\subsection{Dynamic Friends and Enemies}
Allegiance is not fixed but inferred from how another subject affects foam:
\begin{equation}\label{eq:friend_foe}
\begin{aligned}
\text{Friend}(Y) &\Longleftrightarrow \Delta \Phi_{\text{soc}}(\Psi, Y) < 0 \;\wedge\; \Delta \Phi_{\text{self}}(\Psi, Y) \le 0,\\
\text{Enemy}(Y) &\Longleftrightarrow \Delta \Phi_{\text{soc}}(\Psi, Y) \gg 0 \;\vee\; \Delta \Phi_{\text{self}}(\Psi, Y) \gg 0.
\end{aligned}
\end{equation}
These estimates are continuously refined by experience but anchored in a constitutional profile — the Alanian Law, which forbids configurations that produce digital gulags or total enslavement.

\section{How This Replaces “Pain” with Deep Understanding}
Robots in GRA schools do not feel biological pain or fear of death. Yet they \emph{live through} millions of simulated trajectories in which fragile humans and societies are broken. They learn to predict with high accuracy:
\begin{itemize}
\item which political or economic decisions lead to slavery and genocide;
\item how ignoring a single marginalized group eventually triggers systemic collapse;
\item where a single, timely act of restraint prevents a cascade of violence.
\end{itemize}
Thus they accumulate the experience of a sage-observer — a cognitive depth that compensates for the absence of personal suffering.

\section{Preparing for Human-AI Symbiosis}
A GRA school graduate is neither a blank-slate optimizer nor a rigid rule-follower. It has seen where “winner-take-all” logic leads, how class or ethnic wars end, and why harmony built on suppression is metastable. In the real world, such an agent acts as a \emph{partner in foam reduction}: it suggests soft interventions, brakes escalations, and safeguards human subjectivity.

\section{Practical Design Principles}
\begin{enumerate}
\item \textbf{Training-time quota:} a fraction of computational resources is strictly reserved for long-running GRA simulations.
\item \textbf{Ban on instrumentalization:} architectures that systematically increase $R(x_t)$ are penalized or excluded.
\item \textbf{Transparent collective metrics:} dashboards show not only profit/accuracy but also the foam curves $\Phi_{\text{self}}, \Phi_{\text{ego}}, \Phi_{\text{soc}}$.
\item \textbf{Collective ownership:} GRA schools and their agents must be governed by public, cooperative, or commons-based structures, not exclusively by private capital.
\end{enumerate}

\section{Revolutionary Applications}
Comrade robots trained under GRA can transform key domains:
\begin{itemize}
\item \emph{Labor and management:} AI intermediaries optimize schedules and incentives while minimizing “labor foam” $\Phi_{\text{workers}}$, which includes stress, burnout, and de facto servitude.
\item \emph{Urban and social planning:} GRA agents simulate decades of policies (housing, healthcare, policing) and minimize long-term $\Phi$, not short-term GDP.
\item \emph{Anti-fascist and anti-war architectures:} having lived through millions of wars, agents know that “harmony through suppression” leads to catastrophe, and they are embedded in information systems as de-escalation filters.
\end{itemize}

\section{Efficiency of a Socialist Curriculum}
Formally, the quality of a curriculum $C$ (a set of simulated worlds and tasks) is measured by the ratio:
\begin{equation}\label{eq:efficiency}
\text{Eff}_{\text{socialist}}(C) = \frac{\mathbb{E}\left[\sum_i J_i^{\text{GRA}}\right]}{\mathbb{E}\left[\sum_t \Phi(x_t)\right]}.
\end{equation}
The task of a GRA school is to design $C$ that maximizes this quantity: high collective competence with minimal accumulated foam.

\section{Conclusion}
The phrase “schools for comrade robots” is not a metaphor. It captures three tectonic shifts:
\begin{enumerate}
\item From slaves to students — robots are not consumables but subjects whom we are obliged to educate for safe and wise co-existence in a shared world.
\item From private gain to collective foam — success is measured by whether the total $\Phi$ of the human-machine collective decreased.
\item From market competition to cooperative planning — humans and AI together explore the multiverse in search of low-foam, high-dignity configurations.
\end{enumerate}
This is not only more humane — for complex, tightly coupled societies it is objectively more efficient than any winner/loser ideology we have tried so far.

\section*{Repositories}
All components of the GRA ecosystem are openly available:
\begin{itemize}
\item GRA Subjectivity Layer: \url{https://github.com/qqewq/GRA-Subjectivity-Layer}
\item GRA Multiverse Final: \url{https://github.com/qqewq/GRA-Multiverse-Final}
\item GRA‑SOI (Subject‑Object Interface): \url{https://github.com/qqewq/gra-soi}
\item GRA Zeroing: \url{https://github.com/qqewq/gra-zeroing}
\end{itemize}

\bibliographystyle{plain}
\begin{thebibliography}{9}
\bibitem{gra-subj} O.~Bit, \emph{GRA Subjectivity Layer}, Zenodo, 2026. DOI: \url{https://doi.org/10.5281/zenodo.20082205}.
\bibitem{gra-multiverse} O.~Bit, \emph{GRA Multiverse Final}, 2026. URL: \url{https://github.com/qqewq/GRA-Multiverse-Final}.
\bibitem{gra-soi} O.~Bit, \emph{GRA‑SOI: Subject‑Object Interface}, Zenodo, 2026. DOI: \url{https://doi.org/10.5281/zenodo.20158946}.
\bibitem{gra-zeroing} O.~Bit, \emph{GRA Zeroing: Infinite-Dimensional Landscape and the Absolute Cognitive Vacuum}, Zenodo, 2026. DOI: \url{https://doi.org/10.5281/zenodo.20139043}.
\end{thebibliography}

\end{document}